\scanpage{27}\mspage{26}
\deleted{\unsure{Alors} $F$}

Ainsi, les relations d'ordre $\ll$ et $\preceq$ de, $\unsure{\mathcal F|F}$ se déterminent mutuellement, en termes de $\leq$. Peut-on de $=$ exprimer $\leq$ en termes de $\ll$ et de $\preceq$~?
\margin{??}

\uline{NB} On a
\[
 F\leq G\text{ et }F\preceq G\quad\Longrightarrow\quad F=G
\]
\unsure{Il} faudra que ceci soit vrai dans \deleted{\uline{toute}} contrée -- via la condition que si $F\leq G$, $F\ne G$, $\exists\,H$ \inserted{\unsure{fig.|figure|ferm.}} \uline{non vide} $H$ disjointe de $F$, $H\ll G$ (et si on avait $F\preceq G$, $H$ serait disjointe de $G$ par C5, donc \unsure{on|OK} est ramené à ceci. \deleted{$H\ll G$,}
\[
 H\ll G,\quad
 \underset{\inserted{\substack{\text{(i.e. }H\text{ comp. à }G\text{ et}\\H\cap G=\varnothing}}}{H\text{ disjointe de }G}
 \quad\Longrightarrow\quad H=\varnothing_{F}
\]
En effet, $H\leq G'\preceq G$, \deleted{\unsure{$i.e.$ $R_G$ $|H|$}} et $H$ disjointe de $G$ signifie $H$ disjointe de $G'$ \inserted{par C5}, donc $H=\varnothing$ OK~!)

\deleted{\uline{Famille \unsure{intersection|intersections}}}

Qu'en est-il des autres axiomes des contrées~:

\uline{C5} Si $F\preceq G$, $H$, alors $H$ dis. de $F$ ssi disj. de $G$ trivial, car $|F|=|G|$ et disjonction de $H$ \unsure{avec|ou} $H'$ se voit sur les supports par $|H|\cap|H'|=\varnothing$

\uline{C6} \deleted{\unsure{R|S}}Subdivision induite, on la verra tantôt

\uline{C7} Ext. d'une subd. OK on le vérifiera \unsure{tautologiq\(^t\)|topologiq\(^t\)|de tte façon}.

\uline{\unsure{C$8_0$|C$\varphi_0$|C$\varnothing_0$}} trivial
